\documentclass[../../../include/open-logic-section]{subfiles}

\begin{document}

\olfileid{sth}{choice}{intro}

\olsection{પરિચય}

\crefrange{sth:cardinals::chap}{sth:card-arithmetic::chap}માં આપણે
ગણાંકોને ક્રમવાચક સંખ્યાઓ માનીને ગણાંકોનો સિદ્ધાંત
વિકસાવ્યો. આ રીત સુક્રમિતતાની સ્વયંસિદ્ધિ પર આધાર
રાખે છે. સુક્રમિતતા બીજી એક સ્વયંસિદ્ધિ---પસંદગીની
સ્વયંસિદ્ધિ---ને સમકક્ષ છે, અને તેની સ્વીકાર્યતા અંગે
ગંભીર તત્ત્વચિંતનાત્મક ચર્ચા થઈ છે. આ પ્રકરણમાં આપણા
પ્રશ્નો છે: આ સ્વયંસિદ્ધિનો ઉપયોગ કેવી રીતે થાય છે,
અને તેને વાજબી ઠરાવી શકાય?

\end{document}
